% Structured abstract, at most 250 words. Every number is copied from paper/tables/v3_*.tex; keep it in step with cls_results.tex.
\textbf{Background.} Bone suppression helps radiologists find nodules hidden by ribs, and suppressed images have been proposed as
classifier inputs. A suppressor computes its output from the radiograph, so any gain must come from its learned prior; a fair test must
control for network capacity.

\textbf{Purpose.} To test whether the soft-tissue and bone images of a CT-supervised suppressor improve a chest radiograph classifier,
and for which findings.

\textbf{Methods.} A DenseNet-121 classifier was trained on five inputs: the radiograph; the soft-tissue image; radiograph,
soft-tissue and bone images stacked as channels; the same three as streams of a shared encoder (hybrid); and the radiograph in all
three streams (capacity control). Five seeds on NIH ChestX-ray14 (25,596 test images), VinBigData (2,943) and a paediatric pneumonia set
(3,257, five folds) gave 175 runs. A pre-specified plan fixed the tests: paired cluster bootstrap of seed ensembles, Holm correction.

\textbf{Results.} The hybrid did not improve macro-AUROC over the radiograph (NIH $+0.0001$ $[-0.0022,0.0024]$; VinBigData
$+0.0024$ $[-0.0004,0.0053]$; paediatric $-0.0042$ $[-0.0108,0.0021]$) nor over the capacity control or channel stacking. For nodules
and masses the soft-tissue input raised the mean AUROC on NIH ($+0.0089$ $[0.0040,0.0140]$, Holm $p=0.002$), with a similar,
non-significant estimate on VinBigData ($+0.0083$); on NIH, diffuse findings lost ($-0.0032$ to $-0.0052$). Soft-tissue stream maps hit
nodules and masses more often than radiograph maps (interaction $+0.145$ $[0.092,0.199]$). The suppressor's uncertainty did not
improve selective classification.

\textbf{Conclusions.} Bone-suppressed inputs did not make these classifiers more accurate overall; their largest effect on both adult
datasets was a small gain for nodules and masses, significant on NIH only.
